% Chapter Template

\chapter{Introduction} % Main chapter title

\label{Chapter 1} % Change X to a consecutive number; for referencing this chapter elsewhere, use \ref{ChapterX}

\lhead{Introduction } % Change X to a consecutive number; this is for the header on each page - perhaps a shortened title

\section{Introduction}
A Fall-Detection Device basically monitors the accelerometer and gyropscope readings of the person to which it was attached and tries to identify the fall pattern from the statistical data, which it derives from the sensor data. For analysing those, we will use a trained ML model deployed in the cloud . This device is not only capable of finding the fall patterns, but it can also collect the movement data in terms of statistical features of accelerometer and gyroscopic readings for futur usee, which can be further used to train the deployed model after labelling

\section{Motivation}
Many Elderly people will die due to unexpected fall accidents, to make sure we have continuous monitoring of the movement, what if we have a device that analyses the movement patterns, continuously identifies the fall accidents immediately within a fraction of second, alerting the caretakers.  

\section{Contribution of Project}
This project was made into three different sections Frontend, Backend and Hardware. 

The Frontend and hardware sections are looked after by the V.Tejaraam which includes creation of visually appealing dashboard and interactive labeling page along with hardware simulation and connection.

The system architecture and backend server design along with model creation was looked after by D.Adhithya, he design the endpoints, end to end communication and model for identifying the fall pattern, all the connectivity issues, debugging, logs and managing database was covered in the backend server design itself.  
